\section{Proposed Framework}
\label{sec:proposed_framework}

\subsection{Layered Architecture and Total Pipeline}
To overcome the structural rigidity of deterministic intersection controllers, we propose the Semantic Predictive Graph Reinforcement Learning (SPGRL) framework. The SPGRL architecture is constructed as an end-to-end cyber-physical system, composed of seven tightly coupled computational layers: semantic perception, behavioral perception, predictive forecasting, graph representation learning, sustainability optimization, emergency routing, and reinforcement learning control. 

The complete information flow of the SPGRL pipeline is strictly hierarchical. Unconstrained raw traffic video is first ingested and processed via an offline Video Masked Autoencoder (VideoMAE) to extract 768-dimensional temporal-spatial features. These features are mapped into a continuous probability distribution using Multi-Level Density Estimation (MULDE) and calibrated through a Gaussian Mixture Model (GMM) to output a deterministic semantic anomaly metric ($A_s$). Simultaneously, the video undergoes explicit object detection and tracking via YOLO and DeepSORT to extract trajectories, yielding a kinematic behavioral anomaly metric ($A_b$). Concurrently, historical numerical traffic flows are ingested by a Long Short-Term Memory (LSTM) sequence network to compute future trajectory limits ($F_t$) and confidence bounds ($C_f$). The structural topology of the road network is encoded via a Graph Neural Network (GNN) into spatial node embeddings ($G_t$). Real-time telemetry drives the carbon engine to formulate instantaneous emission bounds ($C_t$), while priority metrics trigger the emergency routing engine ($E_t$). All resulting outputs are concatenated into a high-dimensional unified state vector ($Z_t$). This unified manifold is passed to a Multi-Agent Proximal Policy Optimization (MAPPO) architecture, which dictates the signal transition policy. Finally, the decision is intercepted by a deterministic Safety Shield to enforce collision avoidance before physical signal actuation. Fig.~\ref{fig:spgrl_pipeline} illustrates the complete layered architecture of the proposed SPGRL framework.

The seven computational layers execute distinct functions within this pipeline. 
\textbf{Layer 1: Semantic Perception.} This layer identifies non-geometric scene volatility (e.g., severe weather, debris, structural anomalies) by capturing density deviations in the VideoMAE feature space. It allows the framework to inherently understand scene severity without requiring frame-level supervision.

\textbf{Layer 2: Behavioral Anomaly Extraction.} Providing an explicit counterpart to semantic estimation, this layer models micro-kinematic deviations. By extracting exact bounding box velocities, accelerations, and angular divergences, it explicitly tracks hazardous individual vehicular movements.

\textbf{Layer 3: Traffic Prediction.} To shift control from reactive to proactive, the LSTM forecasting layer projects short-term trajectory horizons, providing the agent with future state approximations alongside strict variance-based uncertainty metrics.

\textbf{Layer 4: Graph Coordination.} The spatial layer ensures intersections do not operate in isolation. Message passing along the directed road network topology grants decentralized nodes explicit awareness of incoming shockwaves from adjacent intersections.

\textbf{Layer 5: Carbon Optimization.} This analytical layer explicitly links instantaneous acceleration profiles to empirical greenhouse gas emissions, establishing an active environmental sustainability penalty within the optimization manifold.

\textbf{Layer 6: Emergency Routing.} This priority layer operates outside the standard flow, executing deterministic pathfinding protocols to detect and route critical response vehicles with absolute preemption.

\textbf{Layer 7: RL Control and Safety Enforcement.} The terminal layer performs the optimal phase mapping using MAPPO with centralized training and decentralized execution, bounded by the Safety Shield to ensure all generated neural actions maintain strict physical safety margins.

This layered decomposition is necessary to ensure modularity, scalability, and extreme fault tolerance. By decoupling the sensory modalities, the architecture supports highly distributed execution across edge nodes, preventing singular sensor failures from inducing catastrophic network collapse. Furthermore, the explicit segregation of state variables heavily enhances interpretability, allowing system operators to trace multi-objective optimization decisions directly back to their semantic or behavioral origins.

\subsection{Module Dependency Graph}
The execution integrity of the SPGRL architecture is governed by a strict computational dependency graph mapping the flow of latent information from raw telemetry to physical actuation.

The data dependencies exhibit structural heterogeneity. A sequential semantic dependency requires that VideoMAE compute spatial-temporal volumes prior to MULDE estimating density scores, which in turn precede the GMM calibration of $A_s$. Similarly, a spatial behavioral dependency strictly links YOLO object proposals to DeepSORT association filters to derive $A_b$. Temporal dependencies dictate that historical traffic vectors are processed by the LSTM to produce $F_t$ and $C_f$. Graph representations inherently rely on static road topology matrices propagated through the GNN to derive $G_t$. Furthermore, independent telemetry feeds govern the Carbon Engine ($C_t$) and the Emergency Router ($E_t$). All independent feature extractors must synchronize to form the unified state tensor $\{A_s, A_b, F_t, C_f, G_t, C_t, E_t\} \rightarrow Z_t$. This state constitutes the sole control dependency for the downstream MAPPO policy, the outputs of which are exclusively gated by the Safety Shield prior to signal actuation.

The computational graph properties heavily favor modern high-performance parallel architectures. As a primarily acyclic execution graph, the independent modalities (VideoMAE, YOLO, LSTM, GNN, Carbon, Emergency) evaluate parallelly, establishing rigorous synchronization barriers only at the $Z_t$ construction phase. This decoupled formulation isolates latency bottlenecks to the vision encoders, making the pipeline exceptionally suited for distributed HPC execution and targeted GPU acceleration.

Crucially, the decoupled dependency graph provides profound fault tolerance. The MAPPO actor relies on a structurally static $Z_t$ dimension. In the event of a missing camera feed, the respective vision indices ($A_s$, $A_b$) regress to baseline, forcing the agent to rely gracefully on localized numerical data. A failed prediction disables horizon forecasting without corrupting immediate perception. Graph failures isolate the node to self-observation, while emergency overrides invoke the Safety Shield to deterministically bypass the neural policy entirely. At no point does partial sensory ablation induce full systemic collapse.

The computational dependency complexity bounds for real-time inference execution are strictly deterministic. The inference time complexity scales as follows: VideoMAE scales at $\mathcal{O}(T \times D)$; MULDE executes at $\mathcal{O}(N)$; GMM evaluates at $\mathcal{O}(KD)$; Behavioral tracking operates at $\mathcal{O}(N)$; the LSTM forward pass is bounded by $\mathcal{O}(WH)$; GNN message passing executes at $\mathcal{O}(V+E)$; Emergency routing resolves in $\mathcal{O}(E+V\log V)$; Fusion concatenation occurs at $\mathcal{O}(|Z_t|)$; and the MAPPO actor inference evaluates at $\mathcal{O}(|Z_t||A|)$.
